\section{复随机变量}

在我们前面的讨论中，我们限定随机变量的取值为实数。事实上我们也能自然地把它拓展到复数上，只是这时我们就不能享受实数的线性连续统性质了。（事实上我们在多元随机变量那里早就悄悄地对随机变量的概念进行过推广了，我们甚至可以更进一步定义随机矩阵的概念。）我们仅在这一节中特别强调实数和复数的区别，在其他部分我们仍然默认随机变量就是实的随机变量。

令\(X, Y\)为两个实随机变量，我们定义\(Z=X+iY\)为一个\textbf{复随机变量}。换句话就是：
\[
\forall e \in \Omega, Z(e) = X(e) + iY(e) \in \mathbb{C}
\]

等价地，我们可以说一个复随机变量就是一个\(\Omega\to \mathbb{C}\)的可测函数（其中\(\Omega\)是样本空间）。于是我们可以对复随机变量之间做各种复数运算（加减乘除共轭等），仍然得到复随机变量（或者实随机变量）。

由于复数上没有自然的序了（\(\mathbb{R}^n\)上面也没有），累积分布就不是很好定义了。但是我们还是可以定义实部和虚部的联合分布，作为一个二元实随机变量。

我们仍然可以类似地定义复变量的数值特征。

均值的定义是和实随机变量完全类似的的，\(\E[Z]=\E[X]+i\E[Y]\)。接着我们就可以定义\(k\)阶原点矩为\(\E[Z^k]\)，\(k\)阶中心矩为\(\E[(Z-\E(Z])^k)\)。

很简单地，均值仍然具有线性性。更多地，它还有共轭性质（我们在这里和后面都用星号表示共轭）：
\[
E(Z^*) = \left(E(Z)\right)^*
\]

不过需要小心的是，我们的协方差和方差是半内积的结构，我们最好让它遵守共轭对称性，并且结果为一个实数。

令\(Z,W\)为两个复随机变量，我们定义它们的协方差为：
\[
\Cov (Z,W) = E\left((Z-\E[Z])(W-\E[W])^*\right)
\]
方差就是
\[
\Dev[Z] = \Cov(Z,Z)= E\left(|Z-E(Z)|^2\right)
\]

\begin{attention}
我们在这里的定义让协方差的第一个位置线性，第二个位置共轭线性。二阶矩与它是不同的，\(\Cov(Z,W*)\)也被称为\textbf{伪协方差}，\(\Cov(Z,Z*)\)也被称为\textbf{伪方差}。
\end{attention}

我们仍然可以定义复随机变量的特征函数。事实上，我们在定义特征函数的时候已经偷摸定义了一个复随机变量\(e^{iXt}\)。这里对于复随机变量\(Z=X+iY\)，我们并不采用和实随机变量一样的定义：因为这样可能导致不收敛。我们定义时仍然将实部虚部分开，当作多元随机变量的特征函数处理：
\[
\phi_Z(s, t) = \E\left[e^{i(sX+tY)}\right],\quad s,t\in \mathbb{R}
\]
或者更复数地，令\(\omega =s+it\)：
\[
\phi_Z(\omega) = \E\left[e^{i\Re(\omega^* Z)}\right]
\]

\subsection{复正态分布}
和实的情况一样，复正态分布也绝对是最经典的复分布。\(Z=X+iY\)服从复正态分布也就是它的实部和虚部放在一起服从二维正态分布。这个没有太多可说的，但是我们一般关注的是一类特殊的复正态分布，即\textbf{圆对称复正态分布}。

所谓圆对称就是随机变量的概率密度关于原点旋转对称，即\(p_Z(z)=p_Z(e^{i\phi} z)\)。事实上由于它的圆对称性，我们可以把它的概率密度写成\(p_Z(|z|)\)的形式。或者可以等价地说：\(|Z|\)和\(\arg Z\)是独立的随机变量，并且\(\arg Z\sim U(0,2\pi)\)。也就是每一个圆对称的复随机变量都可以写成\(Z=Re^{i\theta}\)的形式，其中\(\theta \sim U(0,2\pi)\)。

更进一步，我们还可以知道：如果各阶矩存在的话，那它们都是\(0\)。因为对于任意\(\phi\)都有：
\begin{align*}
\E[Z^k] &= \E[(Ze^{i\phi}]^k)\\
&= e^{ik\phi}\E[Z^k]
\end{align*}
所以必须有\(\E[Z^k]=0\)。特别地，二阶矩（就是伪方差），也为零。或者说\(Z\)与\(Z^*\)不相关。做个正交变换，\(Z\)的实部和虚部也不相关。

甚至更进一步，它的实部和虚部是同分布的：\(Z=X+iY\)和\(iZ=-Y+iX\)同分布。更进一步，\(X\)和\(Y\)都是对称分布。

对于二维正态分布来说，这已经足够说明问题了，也就是说明\(X,Y\)独立同分布。我们把圆对称复正态分布记作\(Z\sim CN(0,\sigma^2)\)。

它的概率密度表达非常简洁。\((X,Y)\)的联合概率密度是
\[
f_{XY}(x,y) = \frac{1}{\pi\sigma^2} \exp\left(-\frac{x^2+y^2}{\sigma^2}\right)
\]
换成复的表达就是：
\[
f_Z(z) = = \frac{1}{\pi\sigma^2} \exp\left(-\frac{|z|^2}{\sigma^2}\right)
\]
\begin{attention}
这是对二维Lebesgue测度而言的。我们这里在\(\mathbb{C}\)上直接继承\(\mathbb{R}^2\)上的测度结构。
\end{attention}

圆对称复正态分布\(CN(0,\sigma^2)\)的特征函数就是：
\[
\phi_Z(\omega) = \exp\left(-\frac{\sigma^2}{4}|\omega|^2 \right)
\]